% ===========================================================================
% PART 0 -- the mathematics, from nothing. Assumes no statistics whatsoever.
% Every symbol used later in the paper is defined here and nowhere else.
% ===========================================================================
\section{The mathematics, from zero}
\label{sec:math}

This section assumes no statistics. It builds, in order, every object the rest of
the paper uses. Nothing here is specific to the dataset; §\ref{sec:object} onward
uses only what is defined below. A reader who already knows this material should
still read §\ref{sec:cluster}, §\ref{sec:mde} and §\ref{sec:atten}, because those
three are where papers in this area go wrong.

\subsection{Sets, functions, and the two symbols that do all the work}

A \emph{set} is a collection of distinct things, written $\{a,b,c\}$. We write
$x \in \mathcal{A}$ for ``$x$ is a member of $\mathcal{A}$'', and $|\mathcal{A}|$
for the number of members.

A \emph{function} $f : \mathcal{A} \to \mathcal{B}$ assigns to every member of
$\mathcal{A}$ exactly one member of $\mathcal{B}$.

Two symbols carry most of the load.

\begin{itemize}
\item $\sum_{i=1}^{n} x_i$ means $x_1 + x_2 + \cdots + x_n$. It is an instruction
      to add. Nothing more.
\item $\mathbf{1}\{P\}$ is the \emph{indicator}: it equals $1$ if the statement
      $P$ is true and $0$ if false. It converts a yes/no fact into a number, and
      it is how every ``what fraction of cases'' quantity in this paper is built.
\end{itemize}

\noindent The single most useful identity in the paper follows immediately from
those two: the average of an indicator \emph{is} a proportion.
\begin{equation}
\frac{1}{n}\sum_{i=1}^{n}\mathbf{1}\{P_i\}
  \;=\; \frac{\text{number of } i \text{ with } P_i \text{ true}}{n}.
\label{eq:indavg}
\end{equation}
Every rate reported in this paper -- how often a decision changes, how often a
veto binds -- is an instance of \cref{eq:indavg}. That is why they can all be
treated with one set of tools.

\subsection{Finite probability}

\begin{definition}[Probability space]
A \emph{finite probability space} is a finite set $\Omega$ of \emph{outcomes}
together with a function $\Prob$ assigning to each $\omega \in \Omega$ a number
$\Prob(\omega) \ge 0$, such that $\sum_{\omega \in \Omega} \Prob(\omega) = 1$.
For a subset $E \subseteq \Omega$ (an \emph{event}),
$\Prob(E) := \sum_{\omega \in E}\Prob(\omega)$.
\end{definition}

Two consequences, both immediate from the definition and both used later.

\step{Sub-step 0.2.1 --- probabilities are bounded.}
$E \subseteq \Omega$ implies $\Prob(E) \le \sum_{\omega\in\Omega}\Prob(\omega) = 1$,
and $\Prob(E) \ge 0$ because it is a sum of non-negative terms. So
$0 \le \Prob(E) \le 1$ always.

\step{Sub-step 0.2.2 --- disjoint events add.}
If $E \cap F = \varnothing$ then no outcome is counted twice, so
$\Prob(E \cup F) = \Prob(E) + \Prob(F)$.

\begin{definition}[Random variable]
A \emph{random variable} is a function $X : \Omega \to \R$. It is not random and
it is not a variable; it is a rule that reads off a number from whatever outcome
occurred. We write $\Prob(X = x)$ for $\Prob(\{\omega : X(\omega) = x\})$.
\end{definition}

\subsection{Expectation: the derivation, not the formula}

Suppose we observe $X$ many times and average the results. If outcome $x$ occurs
a fraction $\Prob(X=x)$ of the time, then in $n$ observations it contributes
$x \cdot n\Prob(X=x)$ to the total. Summing over the possible values and dividing
by $n$:
\begin{equation}
\frac{1}{n}\sum_{\text{values } x} x \cdot n\,\Prob(X=x)
  \;=\; \sum_{x} x\,\Prob(X=x).
\end{equation}
The $n$ cancels. The result does not depend on how many times we look. That
limit-of-the-average is the definition.

\begin{definition}[Expectation]
$\displaystyle \E[X] := \sum_{x} x\,\Prob(X = x)$.
\end{definition}

\begin{lemma}[Linearity]
\label{lem:lin}
For random variables $X,Y$ on the same space and constants $a,b$:
$\E[aX + bY] = a\E[X] + b\E[Y]$.
\end{lemma}

\begin{proof}
Write the expectation as a sum over outcomes rather than values, which is
equivalent because grouping outcomes by their value only reorders a finite sum:
$\E[X] = \sum_{\omega} X(\omega)\Prob(\omega)$. Then
\[
\E[aX+bY] = \sum_\omega \big(aX(\omega)+bY(\omega)\big)\Prob(\omega)
 = a\sum_\omega X(\omega)\Prob(\omega) + b\sum_\omega Y(\omega)\Prob(\omega),
\]
which is $a\E[X]+b\E[Y]$. \emph{No independence was used.} This is why linearity
is the one tool that never needs an assumption, and it is used at every later
step where an assumption would be contentious.
\end{proof}

\begin{corollary}[Expectation of an indicator]
\label{cor:indexp}
$\E[\mathbf{1}\{P\}] = 1\cdot\Prob(P) + 0\cdot(1-\Prob(P)) = \Prob(P)$.
\end{corollary}

\Cref{cor:indexp} is the bridge between ``a proportion I measured'' and ``a
probability I want to talk about'', and every rate in this paper crosses it.

\subsection{Variance: why squared deviation, and not something else}

We want a single number for ``how spread out is $X$''. Let $\mu := \E[X]$.

\step{Sub-step 0.4.1 --- the naive attempt fails.}
$\E[X-\mu] = \E[X]-\mu = 0$ by \cref{lem:lin}, always. Average deviation is
identically zero and carries no information: positive and negative deviations
cancel exactly.

\step{Sub-step 0.4.2 --- two repairs are available.}
Remove the sign either by absolute value, $\E|X-\mu|$, or by squaring,
$\E[(X-\mu)^2]$. Both are legitimate measures of spread. Squaring is chosen for
one reason that matters here and one that does not.

\begin{itemize}
\item \textbf{Matters:} squared deviations \emph{add} for independent quantities
      (\cref{lem:varsum}), so the spread of an average has a closed form. Absolute
      deviations do not add.
\item \textbf{Does not matter:} differentiability, which is a convenience for
      optimisation and irrelevant to this paper.
\end{itemize}

\begin{definition}[Variance and standard deviation]
$\Var(X) := \E[(X-\mu)^2]$ and $\operatorname{sd}(X) := \sqrt{\Var(X)}$. The
square root exists because $\Var(X)$ is an average of squares, hence
$\ge 0$. The standard deviation is in the same units as $X$; the variance is in
squared units, which is why intervals are reported with the former.
\end{definition}

\begin{lemma}[Computational form]
\label{lem:varalt}
$\Var(X) = \E[X^2] - (\E[X])^2$.
\end{lemma}
\begin{proof}
$\E[(X-\mu)^2] = \E[X^2 - 2\mu X + \mu^2]
 = \E[X^2] - 2\mu\E[X] + \mu^2 = \E[X^2] - 2\mu^2 + \mu^2 = \E[X^2]-\mu^2$,
using \cref{lem:lin} and the fact that $\mu$ is a constant.
\end{proof}

\begin{corollary}[Variance of an indicator]
\label{cor:indvar}
If $X = \mathbf{1}\{P\}$ with $\Prob(P) = \pi$, then $X^2 = X$ (since
$0^2=0, 1^2=1$), so by \cref{lem:varalt}
$\Var(X) = \pi - \pi^2 = \pi(1-\pi)$.
\end{corollary}

\Cref{cor:indvar} is why every yes/no rate in this paper has a variance we can
write down without estimating anything extra, and why rates near $0$ or $1$ are
measured more precisely than rates near $\tfrac12$.

\begin{lemma}[Scaling]
\label{lem:varscale}
$\Var(aX+b) = a^2\Var(X)$.
\end{lemma}
\begin{proof}
$\E[aX+b] = a\mu+b$, so the deviation is $aX+b-(a\mu+b) = a(X-\mu)$, whose square
is $a^2(X-\mu)^2$. Take expectations and apply \cref{lem:lin}. The additive
constant $b$ vanishes because shifting everything does not change spread.
\end{proof}

\subsection{Independence, covariance, correlation}

\begin{definition}[Independence]
$X$ and $Y$ are \emph{independent} if
$\Prob(X=x \text{ and } Y=y) = \Prob(X=x)\Prob(Y=y)$ for all $x,y$: knowing one
tells you nothing about the other.
\end{definition}

\begin{definition}[Covariance]
$\Cov(X,Y) := \E[(X-\mu_X)(Y-\mu_Y)]$.
\end{definition}

The sign of the product $(X-\mu_X)(Y-\mu_Y)$ is positive when both deviate the
same way and negative when they deviate oppositely, so the average is positive
for co-movement and negative for counter-movement. Note $\Cov(X,X)=\Var(X)$.

\begin{lemma}[Variance of a sum]
\label{lem:varsum}
$\Var(X+Y) = \Var(X)+\Var(Y)+2\Cov(X,Y)$. If $X,Y$ are independent then
$\Cov(X,Y)=0$ and the variances simply add.
\end{lemma}
\begin{proof}
Let $D_X := X-\mu_X$, $D_Y := Y-\mu_Y$. Then
$X+Y-(\mu_X+\mu_Y) = D_X+D_Y$, and
$(D_X+D_Y)^2 = D_X^2 + 2D_XD_Y + D_Y^2$. Take expectations, apply
\cref{lem:lin}, and read off the three terms. For the second claim: under
independence $\E[D_XD_Y] = \E[D_X]\E[D_Y] = 0\cdot 0 = 0$.
\end{proof}

\begin{remark}[The converse is false, and it matters here]
$\Cov(X,Y)=0$ does \emph{not} imply independence. Covariance detects only linear
co-movement. A quantity can be completely determined by another and still have
zero covariance with it. Every ``uncorrelated, therefore unrelated'' inference in
this literature is invalid for this reason, and §\ref{sec:c1} uses a
\emph{non-linear} measure of dependence precisely to avoid it.
\end{remark}

\begin{theorem}[Cauchy--Schwarz]
\label{thm:cs}
$|\Cov(X,Y)| \le \operatorname{sd}(X)\operatorname{sd}(Y)$.
\end{theorem}
\begin{proof}
Write $D_X, D_Y$ as above and let $t \in \R$. Since a square is never negative,
\[
0 \;\le\; \E\big[(D_X + tD_Y)^2\big]
  \;=\; \E[D_X^2] + 2t\,\E[D_XD_Y] + t^2\,\E[D_Y^2]
  \;=\; \Var(X) + 2t\Cov(X,Y) + t^2\Var(Y).
\]
The right-hand side is a quadratic in $t$ that is never negative. A quadratic
$at^2+bt+c$ with $a \ge 0$ is non-negative for all $t$ if and only if its
discriminant $b^2-4ac \le 0$ (otherwise it has two distinct real roots and dips
below zero between them). Here $a=\Var(Y)$, $b=2\Cov(X,Y)$, $c=\Var(X)$, so
\[
4\Cov(X,Y)^2 - 4\Var(X)\Var(Y) \le 0
\;\Longrightarrow\;
\Cov(X,Y)^2 \le \Var(X)\Var(Y).
\]
Take square roots.
\end{proof}

\begin{definition}[Correlation]
If both standard deviations are non-zero,
\[
\Corr(X,Y) \;=\; r \;:=\; \frac{\Cov(X,Y)}{\operatorname{sd}(X)\operatorname{sd}(Y)}.
\]
\end{definition}

\begin{corollary}
$-1 \le r \le 1$, immediately from \cref{thm:cs}.
\end{corollary}

So $r$ is covariance with the units divided out. It is \emph{not} a percentage,
it is \emph{not} a share of anything, and $r=0.5$ is not ``half''. The only
quantity with a share interpretation is $r^2$, and only under a linear model.

\subsection{Samples, estimators, and the thing that is actually random}

We never see $\Prob$. We see $n$ observations $x_1,\dots,x_n$ and must guess
properties of $\Prob$ from them.

\begin{definition}[Estimator]
An \emph{estimator} is any function of the observed data used to guess a fixed
unknown quantity (the \emph{estimand}). We write $\hat\theta$ for an estimator of
$\theta$.
\end{definition}

\begin{definition}[Bias]
$\operatorname{Bias}(\hat\theta) := \E[\hat\theta] - \theta$. An estimator is
\emph{unbiased} if this is zero: it is right \emph{on average across repetitions
of the whole study}, which says nothing about any single study.
\end{definition}

\begin{remark}[What is random]
$\theta$ is a fixed number. $\hat\theta$ varies because the sample varies. Every
statement below about ``the distribution of $\hat\theta$'' is a statement about
what would happen across hypothetical repetitions of the study, not about
$\theta$. This is the single most misread point in applied statistics and it is
the reason §\ref{sec:ci} states explicitly what a confidence interval does
\emph{not} mean.
\end{remark}

\subsection{The sample mean, and the derivation of the standard error}
\label{sec:se}

Let $X_1,\dots,X_n$ be independent, each with mean $\mu$ and variance $\sigma^2$.
Define $\bar X := \frac1n\sum_i X_i$.

\step{Sub-step 0.7.1 --- the mean is unbiased.}
By \cref{lem:lin}, $\E[\bar X] = \frac1n\sum_i \E[X_i] = \frac1n \cdot n\mu = \mu$.

\step{Sub-step 0.7.2 --- the variance of the mean.}
By \cref{lem:varscale} with $a = 1/n$, then \cref{lem:varsum} with independence:
\begin{equation}
\Var(\bar X) \;=\; \frac{1}{n^2}\Var\Big(\sum_i X_i\Big)
  \;=\; \frac{1}{n^2}\sum_i \Var(X_i)
  \;=\; \frac{n\sigma^2}{n^2}
  \;=\; \frac{\sigma^2}{n}.
\label{eq:varmean}
\end{equation}

\step{Sub-step 0.7.3 --- naming it.}
\begin{definition}[Standard error]
$\se(\bar X) := \operatorname{sd}(\bar X) = \sigma/\sqrt{n}$.
\end{definition}

\begin{remark}[The $\sqrt{n}$ is the whole economics of measurement]
Halving the standard error requires \emph{four times} the data. This single fact
determines every sample-size statement in §\ref{sec:mde}, and it is why the
question ``how many prompts do we need'' has an answer that is often unaffordable.
\end{remark}

\begin{remark}[Independence is doing real work in \cref{eq:varmean}]
If the $X_i$ are positively correlated, \cref{lem:varsum} adds covariance terms
and the true variance is \emph{larger} than $\sigma^2/n$. Using
\cref{eq:varmean} anyway understates the uncertainty. §\ref{sec:cluster}
quantifies exactly how much, because this specific error occurred in this work.
\end{remark}

\subsection{Why $n-1$}

To use $\se = \sigma/\sqrt n$ we need $\sigma^2$, which we also do not know. The
obvious estimator $\frac1n\sum_i(X_i-\bar X)^2$ is biased downward, and the
reason is worth the derivation because it exposes what ``degrees of freedom''
means.

\begin{lemma}
\label{lem:bessel}
$\displaystyle \E\Big[\sum_{i=1}^n (X_i-\bar X)^2\Big] = (n-1)\sigma^2.$
\end{lemma}
\begin{proof}
Expand around the true mean $\mu$ by adding and subtracting it:
\[
\sum_i (X_i-\bar X)^2 = \sum_i \big((X_i-\mu)-(\bar X-\mu)\big)^2
= \sum_i (X_i-\mu)^2 - 2(\bar X-\mu)\sum_i (X_i-\mu) + n(\bar X-\mu)^2 .
\]
Now $\sum_i (X_i-\mu) = n(\bar X - \mu)$, so the middle term is
$-2n(\bar X-\mu)^2$ and the last two combine to $-n(\bar X-\mu)^2$:
\[
\sum_i (X_i-\bar X)^2 = \sum_i (X_i-\mu)^2 - n(\bar X-\mu)^2 .
\]
Take expectations. The first term is $n\sigma^2$. For the second,
$\E[(\bar X-\mu)^2] = \Var(\bar X) = \sigma^2/n$ by \cref{eq:varmean}, so it
contributes $n \cdot \sigma^2/n = \sigma^2$. Hence the expectation is
$n\sigma^2 - \sigma^2 = (n-1)\sigma^2$.
\end{proof}

\begin{definition}[Sample variance]
$\displaystyle s^2 := \frac{1}{n-1}\sum_i (X_i-\bar X)^2$, which by
\cref{lem:bessel} satisfies $\E[s^2]=\sigma^2$.
\end{definition}

\step{The interpretation.} Deviations are taken from $\bar X$, which was computed
from the same data and therefore sits closer to the points than $\mu$ does. One
degree of freedom was spent estimating the centre. Dividing by $n-1$ returns it.

\subsection{The central limit theorem, and precisely what it licenses}

\begin{theorem}[Central limit theorem, stated]
\label{thm:clt}
If $X_1,\dots,X_n$ are independent and identically distributed with mean $\mu$
and finite variance $\sigma^2$, then as $n$ grows the distribution of
$(\bar X - \mu)/(\sigma/\sqrt n)$ approaches the standard normal distribution:
$\Prob(Z \le z) \to \Phi(z)$ where
$\Phi(z) = \int_{-\infty}^{z}\frac{1}{\sqrt{2\pi}}e^{-t^2/2}\,dt$.
\end{theorem}

We do not prove this. What matters is what it does and does not say.

\begin{itemize}
\item \textbf{Says:} the \emph{average} becomes normally distributed regardless of
      the shape of the individual observations. This is why a rate built from
      $0/1$ indicators -- about as non-normal as data gets -- still has a normal
      sampling distribution.
\item \textbf{Does not say:} that the data are normal. They are not, and it does
      not matter.
\item \textbf{Fails when:} observations are dependent (§\ref{sec:cluster}), or
      $n$ is small with a skewed distribution, or the quantity is not an average.
      In this paper the third case arises for correlations, which is why
      §\ref{sec:boot} uses resampling instead of \cref{thm:clt} there.
\end{itemize}

\noindent Three quantiles of $\Phi$ are used throughout and are quoted once here:
\begin{equation}
z_{0.975} = 1.9600,\qquad z_{0.80} = 0.8416,\qquad z_{0.99} = 2.3263 .
\label{eq:quantiles}
\end{equation}
$z_q$ is defined by $\Phi(z_q)=q$: the point with $q$ of the distribution below it.

\subsection{Confidence intervals: derivation, and what they do not mean}
\label{sec:ci}

\step{Sub-step 0.10.1.} By \cref{thm:clt}, with probability $0.95$,
\[
-1.96 \;\le\; \frac{\bar X - \mu}{\se} \;\le\; 1.96 .
\]

\step{Sub-step 0.10.2.} Multiply through by $\se$ (positive, so inequalities are
preserved): $-1.96\,\se \le \bar X - \mu \le 1.96\,\se$.

\step{Sub-step 0.10.3.} Subtract $\bar X$ and multiply by $-1$, which
\emph{reverses} both inequalities:
$\bar X - 1.96\,\se \le \mu \le \bar X + 1.96\,\se$.

\begin{definition}[95\% confidence interval]
$\big[\bar X - 1.96\,\se,\; \bar X + 1.96\,\se\big]$.
\end{definition}

\begin{remark}[What it means]
Across hypothetical repetitions of the entire study, $95\%$ of the intervals so
constructed would contain the fixed unknown $\mu$. The randomness is in the
interval, not in $\mu$.
\end{remark}

\begin{remark}[What it does not mean]
It does \emph{not} mean $\Prob(\mu \in [\ell,u]) = 0.95$ for the particular
interval computed: $\mu$ is a constant and either lies in it or does not. It does
\emph{not} mean $95\%$ of future observations fall in it -- that is a prediction
interval and is much wider. And an interval that excludes zero does not mean the
effect is large; §\ref{sec:mde} separates those two questions.
\end{remark}

\subsection{Hypothesis tests, $z$-statistics, and $p$-values}

\begin{definition}[Null hypothesis]
A fully specified statement about the estimand, denoted $H_0$, typically
$\theta = 0$. It is specified in advance so that the distribution of the
estimator \emph{under} it is known.
\end{definition}

\begin{definition}[$z$-statistic]
$z := \hat\theta/\se(\hat\theta)$: the estimate measured in units of its own
uncertainty. A $z$ of $3$ means the observed effect is three standard errors from
zero.
\end{definition}

\begin{definition}[$p$-value]
$p := \Prob(|Z| \ge |z| \mid H_0)$, the probability of seeing an effect at least
this extreme \emph{if the null were true}.
\end{definition}

\begin{remark}[The three standard misreadings, each false]
$p$ is not $\Prob(H_0 \mid \text{data})$ -- it is conditioned the other way round.
$p$ is not the probability the result is a fluke. And $1-p$ is not the
probability the effect is real. The $p$-value is a statement about data under an
assumption, never about the assumption given data.
\end{remark}

\begin{remark}[Why this paper mostly reports $z$ and intervals]
A $p$-value collapses effect size and precision into one number, so a tiny effect
measured precisely and a large effect measured sloppily can produce the same $p$.
The interval keeps them separate, which is why every finding in §\ref{sec:findings}
carries one.
\end{remark}

\subsection{Power, and the minimum detectable effect}
\label{sec:mde}

A null result is uninterpretable without knowing what the design could have seen.

\begin{definition}[Power]
The probability of rejecting $H_0$ when a specific alternative $\theta = \delta$
is true. Conventionally $1-\beta = 0.80$.
\end{definition}

\begin{theorem}[Minimum detectable effect]
\label{thm:mde}
For a two-sided test at level $\alpha$ with power $1-\beta$, the smallest true
effect detectable with that power is
\[
\mathrm{MDE} \;=\; \big(z_{1-\alpha/2} + z_{1-\beta}\big)\cdot \se .
\]
\end{theorem}
\begin{proof}
\emph{Sub-step 1.} We reject when $|\hat\theta| \ge z_{1-\alpha/2}\se$, by
construction of the level-$\alpha$ test.

\emph{Sub-step 2.} Suppose the truth is $\theta = \delta > 0$. Then
$\hat\theta$ is centred at $\delta$ with standard deviation $\se$, so
$(\hat\theta - \delta)/\se$ is standard normal.

\emph{Sub-step 3.} Rejection (ignoring the negligible wrong-tail probability)
requires $\hat\theta \ge z_{1-\alpha/2}\se$, i.e.
\[
\frac{\hat\theta-\delta}{\se} \;\ge\; z_{1-\alpha/2} - \frac{\delta}{\se}.
\]

\emph{Sub-step 4.} That event has probability $1-\beta$ exactly when the
threshold on the right equals $-z_{1-\beta}$, because
$\Prob(Z \ge -z_{1-\beta}) = 1-\beta$ by symmetry of the normal.

\emph{Sub-step 5.} So $z_{1-\alpha/2} - \delta/\se = -z_{1-\beta}$, giving
$\delta = (z_{1-\alpha/2}+z_{1-\beta})\se$.
\end{proof}

\begin{corollary}[The number used throughout]
At $\alpha = 0.05$, $1-\beta = 0.80$: from \cref{eq:quantiles},
$\mathrm{MDE} = (1.9600+0.8416)\,\se = 2.8016\,\se$.
\end{corollary}

\begin{corollary}[Required sample size]
\label{cor:n}
Since $\se = \sigma/\sqrt n$, detecting an effect of size $\delta$ needs
$n = \big(2.8016\,\sigma/\delta\big)^2$.
\end{corollary}

\begin{remark}[Reading a null]
If $\mathrm{MDE} > |\text{the effect the null was supposed to rule out}|$, the
null is \emph{silence}, not evidence of absence. §\ref{sec:findings} reports an
MDE beside every null for this reason.
\end{remark}

\subsection{Clustering, the intraclass correlation, and the design effect}
\label{sec:cluster}

This subsection exists because the error it describes was made in this work and
had to be corrected.

\step{The setting.} Suppose the data have two levels: $P$ independent
\emph{clusters} (here: prompts), and within each cluster $R$ observations (here:
repeated perturbations of the same prompt). Observations in the same cluster are
\emph{not} independent.

\step{Sub-step 0.13.1 --- the model.} Write
$X_{pr} = \mu + a_p + e_{pr}$, where $a_p$ is a cluster effect with variance
$\sigma_b^2$ (``between''), $e_{pr}$ is an observation effect with variance
$\sigma_w^2$ (``within''), and all terms are independent of each other.

\step{Sub-step 0.13.2 --- the cluster mean.} $\bar X_p = \mu + a_p + \bar e_p$,
and by \cref{lem:varsum} and \cref{eq:varmean},
$\Var(\bar X_p) = \sigma_b^2 + \sigma_w^2/R$.

\step{Sub-step 0.13.3 --- the overall mean.} The $P$ cluster means \emph{are}
independent, so \cref{eq:varmean} applies to them:
\begin{equation}
\Var(\hat\pi) \;=\; \frac{1}{P}\Big(\sigma_b^2 + \frac{\sigma_w^2}{R}\Big).
\label{eq:varclust}
\end{equation}

\step{Sub-step 0.13.4 --- what the naive formula would give.} Treating all $PR$
observations as independent with total variance
$\sigma^2 = \sigma_b^2+\sigma_w^2$ gives $\sigma^2/(PR)$.

\begin{definition}[Intraclass correlation]
$\displaystyle \rho := \frac{\sigma_b^2}{\sigma_b^2+\sigma_w^2}$, the share of
total variance that is between clusters. Equivalently, the correlation between
two observations drawn from the same cluster.
\end{definition}

\begin{theorem}[Design effect]
\label{thm:deff}
\[
\mathrm{DEFF}
\;:=\;
\frac{\text{correct variance }\cref{eq:varclust}}{\text{naive variance}}
\;=\; 1+(R-1)\rho .
\]
\end{theorem}
\begin{proof}
Divide \cref{eq:varclust} by $\sigma^2/(PR)$:
\[
\frac{\frac{1}{P}(\sigma_b^2+\sigma_w^2/R)}{\frac{\sigma_b^2+\sigma_w^2}{PR}}
= \frac{R\sigma_b^2+\sigma_w^2}{\sigma_b^2+\sigma_w^2}
= \frac{R\sigma_b^2+\sigma_w^2 + (\sigma_b^2-\sigma_b^2)}{\sigma_b^2+\sigma_w^2}
= 1 + \frac{(R-1)\sigma_b^2}{\sigma_b^2+\sigma_w^2}
= 1+(R-1)\rho. \qedhere
\]
\end{proof}

\begin{corollary}
Standard errors computed as if the $PR$ observations were independent are too
small by a factor $\sqrt{1+(R-1)\rho}$, and $z$-statistics are too large by the
same factor.
\end{corollary}

\begin{remark}[Effective sample size]
$n_{\mathrm{eff}} = PR/\mathrm{DEFF}$. With $R=20$ and $\rho = 0.27$,
$\mathrm{DEFF}=6.13$, so $19{,}360$ rows carry the information of about
$3{,}150$ -- and the honest count of independent units remains $P = 968$.
\end{remark}

\subsection{The bootstrap}
\label{sec:boot}

The formulas above give $\se$ for a \emph{mean}. For a correlation, a median, a
ratio, or anything else, no such formula is available in general. The bootstrap
replaces it with computation.

\begin{protocol}[Cluster bootstrap]
\label{prot:boot}
Given clusters $c_1,\dots,c_P$ and a statistic $T$:
\begin{enumerate}
\item Draw $P$ clusters \emph{with replacement} from $\{c_1,\dots,c_P\}$.
\item Compute $T$ on the resampled collection; call it $T^{(1)}$.
\item Repeat $B$ times, obtaining $T^{(1)},\dots,T^{(B)}$.
\item Report the $2.5$th and $97.5$th percentiles as a $95\%$ interval.
\end{enumerate}
\end{protocol}

\step{Why it works, in one line.} The variation of $T$ across resamples of the
observed data mimics the variation of $T$ across fresh samples from the
population, because the observed sample is our best available stand-in for that
population.

\step{Why the resampling unit must be the cluster.} Resampling \emph{rows}
would create datasets whose within-cluster dependence differs from the real one,
and would reproduce the understatement of \cref{thm:deff}. Every interval in
this paper resamples prompts.

\step{When it fails.} With few clusters, at the boundary of a parameter's range,
or for statistics that are not smooth functions of the data. All three are
checked in §\ref{sec:limits}.

\subsection{Permutation tests, and exactly what they test}

\begin{protocol}[Permutation test]
To test whether the \emph{pairing} between two vectors $x$ and $y$ carries
information: compute $T(x,y)$; then repeatedly shuffle $x$ and recompute
$T(x_{\text{shuffled}}, y)$; the $p$-value is the fraction of shuffles reaching
$|T|$ at least as large as observed.
\end{protocol}

\begin{remark}[The exchangeability assumption]
Shuffling assumes that under $H_0$ all orderings are equally likely. If $x$ has
internal structure -- clustering, time order -- shuffling destroys that too, and
the test then answers a question nobody asked. In §\ref{sec:findings} the shuffle
is applied to cluster-level values for this reason.
\end{remark}

\begin{remark}[Resolution floor]
With $B$ shuffles the smallest attainable $p$ is $1/(B+1)$: a reported
$p = 0.0000$ from $B = 200$ means only $p < 0.005$, and no more.
\end{remark}

\begin{remark}[What a permutation null cannot do]
It answers \emph{did the pairing matter}, never \emph{why}. It is a falsifier,
never a confirmer of a mechanism.
\end{remark}

\subsection{Multiplicity}

\begin{theorem}[Bonferroni]
\label{thm:bonf}
If $C$ tests are each performed at level $\alpha/C$, then the probability of at
least one false rejection when all nulls are true is at most $\alpha$.
\end{theorem}
\begin{proof}
Let $A_i$ be the event of a false rejection on test $i$, so
$\Prob(A_i)=\alpha/C$. The union bound gives
$\Prob(\bigcup_i A_i) \le \sum_i \Prob(A_i) = C\cdot\alpha/C = \alpha$. The union
bound itself follows from \cref{cor:indexp} plus
$\mathbf{1}\{\bigcup A_i\} \le \sum_i \mathbf{1}\{A_i\}$, which holds pointwise
because the left side is $0$ or $1$ and the right side counts.
\end{proof}

\begin{corollary}[The bar used in this paper]
For $C \approx 170$ tests at $\alpha = 0.05$: $\alpha/C \approx 0.0003$, whose
two-sided normal quantile is $|z| \approx 3.6$. This paper uses the more
conservative $|z| > 3.9$ and reports the number of tests and the number of
survivors together, since reporting only survivors is the multiplicity failure
with manners.
\end{corollary}

\begin{remark}[Bonferroni versus false discovery rate]
Bonferroni controls the probability of \emph{any} false positive, which is the
right target when a single false claim is costly. The
Benjamini--Hochberg procedure instead controls the expected \emph{share} of false
positives among rejections, and rejects hypothesis $(k)$ when
$p_{(k)} \le qk/C$ for the largest such $k$. Its threshold at rank $k=C$ is $q$
itself, not $q/C$ -- confusing the two demands orders of magnitude more
resolution than BH actually requires.
\end{remark}

\subsection{Measurement error and attenuation}
\label{sec:atten}

Every quantity in §\ref{sec:object} is measured by a fallible instrument. This
subsection derives the exact consequence.

\step{Sub-step 0.17.1 --- the model.} Suppose the truth is $T$ but we observe
$X = T + \eps$, with $\eps$ independent of $T$, $\E[\eps]=0$,
$\Var(\eps)=\sigma_\eps^2$.

\step{Sub-step 0.17.2 --- reliability.} Define
\[
\mathrm{rel} \;:=\; \frac{\Var(T)}{\Var(X)} = \frac{\sigma_T^2}{\sigma_T^2+\sigma_\eps^2}
\;\in\;[0,1],
\]
the share of observed variance that is signal. \Cref{lem:varsum} with
independence gives the denominator.

\step{Sub-step 0.17.3 --- two independent measurements.} If $X_1 = T+\eps_1$ and
$X_2 = T+\eps_2$ with $\eps_1 \perp \eps_2$ and equal error variances, then
$\Cov(X_1,X_2) = \Var(T) = \sigma_T^2$ (the error terms contribute nothing), and
each has variance $\sigma_T^2+\sigma_\eps^2$. Therefore
\begin{equation}
\Corr(X_1,X_2) \;=\; \frac{\sigma_T^2}{\sigma_T^2+\sigma_\eps^2} \;=\; \mathrm{rel}.
\label{eq:relcorr}
\end{equation}
\emph{The correlation between two independent measurements of the same thing
\underline{is} the reliability of one of them.} This is what makes
§\ref{sec:instrument} computable at zero extra cost.

\begin{theorem}[Attenuation]
\label{thm:atten}
Let $X = T+\eps_X$ and $Y = U+\eps_Y$ with independent errors. Then
\[
\Corr(X,Y) \;=\; \Corr(T,U)\cdot\sqrt{\mathrm{rel}_X\,\mathrm{rel}_Y}.
\]
\end{theorem}
\begin{proof}
$\Cov(X,Y) = \Cov(T+\eps_X, U+\eps_Y) = \Cov(T,U)$, since all three cross terms
involve an error independent of everything else. The denominators are
$\operatorname{sd}(X) = \sigma_T/\sqrt{\mathrm{rel}_X}$ and likewise for $Y$,
because $\mathrm{rel}_X = \sigma_T^2/\sigma_X^2$ rearranges to
$\sigma_X = \sigma_T/\sqrt{\mathrm{rel}_X}$. Substituting,
\[
\Corr(X,Y) = \frac{\Cov(T,U)}{\frac{\sigma_T}{\sqrt{\mathrm{rel}_X}}\cdot\frac{\sigma_U}{\sqrt{\mathrm{rel}_Y}}}
= \frac{\Cov(T,U)}{\sigma_T\sigma_U}\sqrt{\mathrm{rel}_X\mathrm{rel}_Y}. \qedhere
\]
\end{proof}

\begin{corollary}[The ceiling]
A perfect true relationship, $\Corr(T,U)=1$, reads at most
$\sqrt{\mathrm{rel}_X \mathrm{rel}_Y}$. No amount of data raises this: it is
bias, not noise, and $\sqrt{n}$ does not touch it.
\end{corollary}

\subsection{Reliability of an aggregate, and the Spearman--Brown formula}

Averaging several noisy measurements reduces error. The exact statement:

\begin{theorem}[Spearman--Brown]
\label{thm:sb}
If a score is the sum of $m$ exchangeable components each with reliability
correlation $r$ to another such component, then doubling the number of components
gives reliability $2r/(1+r)$.
\end{theorem}
\begin{proof}[Derivation]
Let each component be $T + \eps_i$ with $\Var(T)=\sigma_T^2$,
$\Var(\eps_i)=\sigma_\eps^2$, errors independent. The sum of $m$ components is
$mT + \sum_i \eps_i$, with signal variance $m^2\sigma_T^2$ and error variance
$m\sigma_\eps^2$ by \cref{lem:varsum}. So
\[
\mathrm{rel}_m = \frac{m^2\sigma_T^2}{m^2\sigma_T^2+m\sigma_\eps^2}
= \frac{m\sigma_T^2}{m\sigma_T^2+\sigma_\eps^2}
= \frac{m\lambda}{m\lambda+1},\qquad \lambda := \sigma_T^2/\sigma_\eps^2 .
\]
For $m=1$ this is $r = \lambda/(\lambda+1)$, so $\lambda = r/(1-r)$. Substituting
$m=2$:
\[
\mathrm{rel}_2 = \frac{2r/(1-r)}{2r/(1-r)+1} = \frac{2r}{2r+(1-r)} = \frac{2r}{1+r}. \qedhere
\]
\end{proof}

\begin{remark}[Use here]
Splitting the annotators of a prompt into two halves and correlating the two
half-scores estimates the reliability of a \emph{half}-sized panel. Applying
\cref{thm:sb} corrects it to the full panel, which is the quantity that actually
enters \cref{thm:atten}.
\end{remark}

\subsection{Two-way clustering}

When observations are grouped along two crossed dimensions at once -- here,
prompts and annotators, since an annotator appears on several prompts and a
prompt is rated by several annotators -- neither one-way correction suffices.

\begin{theorem}[Cameron--Gelbach--Miller]
\label{thm:cgm}
$\hat V_{\text{two-way}} = \hat V_{\text{prompt}} + \hat V_{\text{annotator}}
 - \hat V_{\text{prompt}\times\text{annotator}}$.
\end{theorem}
\begin{proof}[Sketch]
Two observations are dependent if they share a prompt \emph{or} an annotator. By
inclusion--exclusion applied to the covariance terms of \cref{lem:varsum}, pairs
sharing \emph{both} are counted once in each of the first two sums and must be
subtracted once. The estimator can be slightly negative in finite samples, which
signals that one dimension carries almost no dependence.
\end{proof}

\subsection{Summary of Part 0}

Everything the rest of the paper needs is now defined: expectation
(\cref{lem:lin}), variance (\cref{lem:varalt}, \cref{cor:indvar}), the standard
error (\cref{eq:varmean}), correlation and its bound (\cref{thm:cs}), intervals
(§\ref{sec:ci}), power (\cref{thm:mde}), clustering (\cref{thm:deff}), resampling
(\cref{prot:boot}), multiplicity (\cref{thm:bonf}), attenuation
(\cref{thm:atten}) and reliability (\cref{thm:sb}). No further statistical
machinery is introduced.
